\label{chap:cantor}

Throughout this chapter $b \geq 3$, $D \subset \{0,\dots,b-1\}$ with $0 \in D$ and
$2 \leq \abs{D} < b$, $\delta \coloneqq \abs{D}$, and
$\Cantor = \Cantor_{b,D} = \{0 = k_0 < k_1 < \cdots\}$.

Everything in this chapter is elementary and finitary.  It is the recommended first stage of
the formalization.

\section{Definitions}

\begin{definition}[Integer Cantor set]
  \label{def:cantor}
  \lean{CantorSzemeredi.CantorData, CantorSzemeredi.CantorData.enum, CantorSzemeredi.CantorData.set}
  \leanok
  Given $b$ and $D$ as above, set
  $\Cantor_{b,D} \coloneqq \{ \sum_{j=0}^{k} d_j b^j : k \in \NNz,\ d_j \in D \}$,
  and let $k_n$ denote the $n$-th element in increasing order, $k_0 = 0$.
\end{definition}

\begin{remark}[Lean design note]
  In Lean it is better to \emph{define} the enumeration by the formula of
  Lemma~\ref{lem:diff}(i) --- that is, as the digit-transfer map $\varphi$ from base $\delta$
  to base $b$ --- and then \emph{prove} that $\varphi$ is a strictly monotone bijection onto
  $\Cantor_{b,D}$.  This makes Lemma~\ref{lem:diff}(i) true by definition and reduces
  Lemma~\ref{lem:diff}(ii) to a digit computation, rather than forcing a development of
  ``the $n$-th element of a set of naturals in increasing order''.  Mathlib's
  \texttt{Nat.digits} and \texttt{Nat.ofDigits} are the right tools.
\end{remark}

\begin{definition}[Syndetic set]
  \label{def:syndetic}
  \lean{CantorSzemeredi.Syndetic}
  \leanok
  A set $B \subset \NNz$ is \emph{syndetic} if there is an $L$ with
  $B \cap [n, n+L] \neq \emptyset$ for every $n$.
\end{definition}

\begin{definition}[Upper Banach density]
  \label{def:banach-density}
  \lean{CantorSzemeredi.upperBanachDensity}
  \leanok
  For $A \subset \ZZ$,
  $\dstar(A) \coloneqq \limsup_{N - M \to \infty} \abs{A \cap [M,N)} / (N-M)$.
\end{definition}

\begin{definition}[Uniform average and uniform liminf]
  \label{def:uniform-liminf}
  \lean{CantorSzemeredi.uniformAvg, CantorSzemeredi.uniformLiminf}
  \leanok
  For a bounded sequence $(a_n)_{n \geq 0}$ put $\unifavg a_n = \frac{1}{N-M}\sum_{M \leq n < N} a_n$,
  and $\unifliminf \unifavg a_n \coloneqq \sup_{L} \inf \{ \avg{M \leq n < M+L'} a_n : M \geq 0,\ L' \geq L \}$.
\end{definition}

\begin{remark}[Lean design note]
  Definition~\ref{def:uniform-liminf} should be set up once, as a filter: the uniform limit is
  the limit along the filter on pairs $(M,N)$ with $M < N$ generated by the sets
  $\{(M,N) : N - M \geq L\}$.  Every ``$\uniflim$'' and ``$\unifliminf$'' in this blueprint is
  then a \texttt{Filter.Tendsto} / \texttt{Filter.liminf} statement along that single filter,
  and the many routine manipulations (a uniform limit dominates a limit along any one sequence
  of intervals of growing length, etc.) are proved once.
\end{remark}

\begin{lemma}[Positive uniform liminf gives a syndetic set]
  \label{lem:pos-liminf-syndetic}
  \lean{CantorSzemeredi.syndetic_of_uniformLiminf_pos}
  \leanok
  \uses{def:syndetic, def:uniform-liminf}
  Let $(a_n)_{n \geq 0}$ be non-negative with $\unifliminf \unifavg a_n > 0$.  Then
  $\{ n : a_n > 0 \}$ is syndetic.
\end{lemma}

\begin{proof}
  \uses{def:syndetic, def:uniform-liminf}
  \leanok
  Write $c > 0$ for the uniform liminf.  If the set were not syndetic there would be intervals
  $[M_i, N_i)$ with $N_i - M_i \to \infty$ containing no such index; the average over each is
  then $0$, contradicting $c > 0$.
\end{proof}

\section{The block structure of the enumeration}

\begin{lemma}[Enumeration of the Cantor set]
  \label{lem:diff}
  \lean{CantorSzemeredi.CantorData.enum_strictMono, CantorSzemeredi.CantorData.mem_set_iff,
    CantorSzemeredi.CantorData.enum_block}
  \leanok
  \uses{def:cantor}
  Let $D = \{0 = d_0 < d_1 < \cdots < d_{\delta-1}\}$.  Then:
  \begin{enumerate}
    \item[(i)] For $n \in \NNz$ written in base $\delta$ as $n = \sum_{c=0}^{t} n_c \delta^c$
      with $n_c \in \{0,\dots,\delta-1\}$, we have $k_n = \sum_{c=0}^{t} d_{n_c} b^c$.
    \item[(ii)] For $i \geq 0$, $n \geq 0$ and $0 \leq j < \delta^i$,
      \[ k_{\delta^i n + j} = b^i k_n + k_j . \]
  \end{enumerate}
\end{lemma}

\begin{proof}
  \uses{def:cantor}
  \leanok
  For (i), define $\varphi(n) \coloneqq \sum_c d_{n_c} b^c$.  Since $d_0 = 0$, leading zeros do
  not change $\varphi(n)$, and the base $b$ digits of $\varphi(n)$ are exactly the $d_{n_c}$.
  Hence $\varphi$ maps $\NNz$ onto $\Cantor$.  It is strictly increasing --- compare at the
  highest position where the base $\delta$ digits differ, using that $c \mapsto d_c$ is
  increasing --- hence a strictly monotone bijection $\NNz \to \Cantor$, so $\varphi(n) = k_n$.

  For (ii), the base $\delta$ digits of $\delta^i n + j$ are those of $j$ in positions
  $0,\dots,i-1$ and those of $n$ shifted up by $i$; apply (i).
\end{proof}

\begin{corollary}[Self-similarity]
  \label{cor:self-similar}
  \lean{CantorSzemeredi.CantorData.pow_mul_mem, CantorSzemeredi.CantorData.enum_pow_mul}
  \leanok
  \uses{lem:diff}
  For $i \geq 0$, $b^i \Cantor \subset \Cantor$; explicitly $b^i k_n = k_{\delta^i n}$.
\end{corollary}

\begin{proof}
  \uses{lem:diff}
  \leanok
  Lemma~\ref{lem:diff}(ii) with $j = 0$, using $k_0 = 0$.
\end{proof}

\begin{corollary}[Absorption of small translates]
  \label{cor:absorb}
  \lean{CantorSzemeredi.CantorData.enum_add_mem_block,
    CantorSzemeredi.CantorData.pow_mul_enum_add_mem}
  \leanok
  \uses{lem:diff}
  Let $i \geq 0$, let $a \in \Cantor \cap [0, b^i)$ and write $\iota(a) < \delta^i$ for the
  index of $a$, so $a = k_{\iota(a)}$.  Then for every $u \geq 0$,
  \begin{equation}
    \label{eq:absorb}
    b^i k_u + a = k_{\delta^i u + \iota(a)} \in \Cantor ,
  \end{equation}
  and the index $\delta^i u + \iota(a)$ lies in the block $[\delta^i u, \delta^i(u+1))$.
\end{corollary}

\begin{proof}
  \uses{lem:diff}
  \leanok
  Immediate from Lemma~\ref{lem:diff}(ii).
\end{proof}

\begin{remark}
  Corollary~\ref{cor:absorb} is exactly what makes the Furstenberg--Katznelson--Ornstein
  compact-extension argument (which produces return times $n + a$, a \emph{translate}) usable
  along a Cantor set, whereas the van der Waerden argument (which produces a \emph{dilate} $sn$)
  is not.  A Cantor set is not closed under dilation.
\end{remark}

\begin{lemma}[Differences along a shift]
  \label{lem:differences}
  \lean{CantorSzemeredi.CantorData.enum_sub_enum_eq,
    CantorSzemeredi.CantorData.card_goodResidues_ge}
  \leanok
  \uses{lem:diff}
  Fix integers $h \geq 1$ and $\ell \geq 1$ with $\delta^\ell > h$.  For
  $0 \leq c < \delta^\ell - h$ set $P_c \coloneqq \{ n \geq 0 : n \equiv c \bmod \delta^\ell \}$
  and $\mathbb{P} \coloneqq \bigcup_{0 \leq c < \delta^\ell - h} P_c$.  Then
  \begin{enumerate}
    \item[(i)] for every $n \in P_c$, $k_{n+h} - k_n = \Delta_c \coloneqq k_{c+h} - k_c$;
    \item[(ii)] $\abs{\mathbb{P} \cap [M,N)} \geq (1 - h\delta^{-\ell})(N-M) - \delta^\ell$ for every
      interval $[M,N)$.
  \end{enumerate}
  So on a set of indices of uniform density at least $1 - h\delta^{-\ell}$ the difference
  $k_{n+h} - k_n$ takes at most $\delta^\ell$ values, each on a union of arithmetic progressions
  of step $\delta^\ell$.
\end{lemma}

\begin{proof}
  \uses{lem:diff}
  \leanok
  For (i), write $n = \delta^\ell a + c$ with $0 \leq c < \delta^\ell - h$; then also
  $n + h = \delta^\ell a + (c+h)$ with $0 \leq c+h < \delta^\ell$, so Lemma~\ref{lem:diff}(ii)
  gives $k_n = b^\ell k_a + k_c$ and $k_{n+h} = b^\ell k_a + k_{c+h}$.  Subtract.
  For (ii), $\mathbb{P}$ is a union of $\delta^\ell - h$ residue classes mod $\delta^\ell$.
\end{proof}

\section{Residues}

\begin{lemma}[Syndetic residues]
  \label{lem:syndetic-residues}
  \lean{CantorSzemeredi.CantorData.syndetic_dvd_enum}
  \leanok
  \uses{lem:diff, def:syndetic}
  For every integer $q \geq 1$ the set $S_q \coloneqq \{ n \in \NNz : q \mid k_n \}$ is syndetic.
\end{lemma}

\begin{proof}
  \uses{lem:diff, def:syndetic}
  \leanok
  Work in $\ZZ/q$ and set, for $i \geq 0$,
  \[
    R_i \coloneqq \{ k_j \bmod q : 0 \leq j < \delta^i \}, \qquad
    A_c \coloneqq \{ d b^c \bmod q : d \in D \} .
  \]
  Splitting $0 \leq j < \delta^{i+1}$ as $j = \delta^i c + j'$ with $c < \delta$ and
  $j' < \delta^i$, Lemma~\ref{lem:diff}(ii) and $k_c = d_c$ give
  $k_j = b^i d_c + k_{j'}$, that is
  \begin{equation}
    \label{eq:Rsucc}
    R_{i+1} = R_i + A_i .
  \end{equation}
  Since $0 \in D$ we have $0 \in A_i$, so the $R_i$ increase; being subsets of the finite set
  $\ZZ/q$ they stabilize: there is $i_0$ with $R_i = R_{i_0} \eqqcolon W$ for all $i \geq i_0$.
  By \eqref{eq:Rsucc}, $W + A_i = W$ for every $i \geq i_0$.

  Let $H \coloneqq \{ t \in \ZZ/q : W + t = W \}$, a subgroup of $\ZZ/q$.  For $i \geq i_0$ and
  $a \in A_i$ the translate $W + a$ is contained in $W + A_i = W$ and has the same cardinality
  as $W$, so $W + a = W$; thus $A_i \subset H$ for every $i \geq i_0$.  Since $0 \in W$ we also
  have $H \subset W$.

  Now fix $i \geq i_0$.  Expanding $k_n$ along the recursion $k_n = d_{n \bmod \delta} + b
  k_{\lfloor n/\delta \rfloor}$ and inducting on $n$ shows $b^i k_n \bmod q \in H$ for every
  $n$: each term is an element of some $A_{i'}$ with $i' \geq i$.  Hence
  $-b^i k_n \bmod q \in H \subset W = R_i$, so there is $0 \leq j < \delta^i$ with
  $k_j \equiv -b^i k_n$, and then $k_{\delta^i n + j} = b^i k_n + k_j \equiv 0 \bmod q$ by
  Lemma~\ref{lem:diff}(ii).  So $S_q$ meets every block $[\delta^i n, \delta^i(n+1))$, and is
  syndetic with gaps at most $2 \delta^i$.

  (This replaces the argument via eventual periodicity of $b^c \bmod q$ given in the paper:
  the stabilization of the $R_i$ already supplies everything, and no distinct-exponent
  realization of sums is needed.)
\end{proof}

\begin{remark}
  Lemma~\ref{lem:syndetic-residues} is used only for
  Theorem~\ref{thm:compact-roots-of-unity}, which is not needed for the main theorem.
\end{remark}

\section{A uniform exponential sum estimate}

Write $e(\theta) \coloneqq \eee^{2\pi i \theta}$ and $\norm{\theta}$ for the distance from
$\theta$ to the nearest integer.

\begin{definition}
  \label{def:Phi}
  \lean{CantorSzemeredi.CantorData.Phi}
  \leanok
  \uses{def:cantor}
  For $\beta \in \RR$ and $L \geq 1$,
  $\Phi_L(\beta) \coloneqq \sup_{M \geq 0} \bigl\lvert \avg{M \leq n < M+L} e(\beta k_n) \bigr\rvert$.
\end{definition}

\begin{lemma}[Block factorization, \eqref{eq:Pell}]
  \label{lem:blockAvg}
  \lean{CantorSzemeredi.CantorData.blockAvg, CantorSzemeredi.CantorData.blockAvg_eq_prod,
    CantorSzemeredi.CantorData.norm_blockAvg_le_one}
  \leanok
  \uses{lem:diff}
  Let $\rho(\theta) \coloneqq \avg{d \in D} e(d\theta)$ and
  $P_\ell(\beta) \coloneqq \avg{0 \leq j < \delta^\ell} e(\beta k_j)$.  Then
  \[ P_\ell(\beta) = \prod_{c=0}^{\ell-1} \rho(b^c \beta) , \]
  and in particular $\abs{P_\ell(\beta)} \leq 1$.
\end{lemma}

\begin{proof}
  \leanok
  \uses{lem:diff}
  By Lemma~\ref{lem:diff}(ii), peeling the top base $\delta$ digit of the index factors the sum
  over a block of length $\delta^{\ell+1}$ as a digit sum at scale $b^\ell$ times the sum over a
  block of length $\delta^\ell$; induct on $\ell$.
\end{proof}

\begin{lemma}[Window split, \eqref{eq:window-split}]
  \label{lem:window-split}
  \lean{CantorSzemeredi.CantorData.norm_sum_Ico_le, CantorSzemeredi.CantorData.Phi_le}
  \leanok
  \uses{lem:diff, lem:blockAvg, def:Phi}
  For every window $[M,N)$ of indices and every $i \geq 0$,
  \[
    \Bigl\lvert \sum_{M \leq n < N} e(\beta k_n) \Bigr\rvert
      \leq (N-M)\abs{P_i(\beta)} + 2\delta^i ,
  \]
  and consequently $\Phi_L(\beta) \leq \abs{P_i(\beta)} + 2\delta^i/L$.  Neither bound depends
  on $M$.
\end{lemma}

\begin{proof}
  \leanok
  \uses{lem:diff, lem:blockAvg}
  If $[M,N)$ contains a full block, split it into the indices below $\delta^i\lceil M/\delta^i\rceil$,
  the full blocks, and the indices above $\delta^i\lfloor N/\delta^i\rfloor$; there are fewer than
  $\delta^i$ of the former and of the latter, and by Lemma~\ref{lem:diff}(ii) the sum over each
  full block is $\delta^i P_i(\beta) e(\beta b^i k_m)$, of modulus at most $\delta^i\abs{P_i(\beta)}$.
  Otherwise $N - M < 2\delta^i$ and the trivial bound suffices.
\end{proof}

\begin{lemma}[Uniform exponential sum estimate]
  \label{lem:irrational}
  \lean{CantorSzemeredi.CantorData.excSet, CantorSzemeredi.CantorData.norm_rhoRed_lt_one,
    CantorSzemeredi.CantorData.exists_kappa, CantorSzemeredi.exists_nrm_ge,
    CantorSzemeredi.CantorData.exists_norm_blockAvg_lt,
    CantorSzemeredi.CantorData.tendsto_Phi_zero,
    CantorSzemeredi.CantorData.countable_excSet}
  \leanok
  \uses{lem:diff, def:Phi, lem:window-split, lem:blockAvg}
  Let $E_0 \subset \TT$ be the set of $\beta$ for which $\gcd(D)\,\beta$ is a $b$-adic rational
  modulo $1$; it is countable.  Then for every $\beta \notin E_0$,
  \[ \lim_{L \to \infty} \Phi_L(\beta) = 0 . \]
\end{lemma}

\begin{remark}[Lean design note]
  The formalization does not build the auxiliary Cantor set $\Cantor_{b,D'}$ of the reduction
  below.  Instead it reduces only the \emph{digit average}: with $s = \gcd(D)$ it defines
  $\rho'(\theta) \coloneqq \avg{d \in D} e((d/s)\theta)$
  (\texttt{CantorData.rhoRed}), so that $\rho(\theta) = \rho'(s\theta)$ and the reduced
  digits have gcd $1$.  Everything after the reduction is then carried out for $\rho'$, and
  the enumeration $k_n$ is never changed.
\end{remark}

\begin{proof}
  \leanok
  \uses{lem:diff, def:Phi}
  \emph{Reduction to $\gcd(D) = 1$.}  With $s \coloneqq \gcd(D)$ and
  $D' \coloneqq \{d/s : d \in D\}$ we have $k_n = s k'_n$ and $e(\beta k_n) = e(s\beta\, k'_n)$,
  and $s\beta$ is $b$-adic rational iff $\beta \in E_0$.  So assume $s = 1$; then $E_0$ is the
  set of $b$-adic rationals mod $1$.

  \emph{Digit average.}  Let $\rho(\theta) \coloneqq \avg{d \in D} e(d\theta)$.  Then
  $\abs{\rho(\theta)} \leq 1$ with equality only if $\theta \in \ZZ$: equality forces all
  $e(d\theta)$ to be equal, hence $=1$ since $0 \in D$, so $d\theta \in \ZZ$ for all $d \in D$
  and then $\theta \in \ZZ$ as $\gcd(D) = 1$.

  \emph{Products over blocks.}  By Lemma~\ref{lem:diff}(i), as $j$ runs over $0 \leq j < \delta^\ell$
  the base $b$ digits of $k_j$ in positions $0,\dots,\ell-1$ run independently over $D$, so
  \begin{equation}
    \label{eq:Pell}
    P_\ell(\beta) \coloneqq \avg{0 \leq j < \delta^\ell} e(\beta k_j) = \prod_{c=0}^{\ell-1} \rho(b^c \beta) .
  \end{equation}

  \emph{Window split.}  Fix $i \geq 1$ and $[M,N)$, and set $M_i = \lceil M/\delta^i \rceil$,
  $N_i = \lfloor N/\delta^i \rfloor$.  At most $2\delta^i$ indices of $[M,N)$ lie outside
  $[\delta^i M_i, \delta^i N_i)$, and the rest are partitioned into blocks
  $\{\delta^i m + j : 0 \leq j < \delta^i\}$.  By Lemma~\ref{lem:diff}(ii) the sum over such a
  block is $\delta^i P_i(\beta) e(\beta b^i k_m)$, of modulus $\leq \delta^i \abs{P_i(\beta)}$.
  Hence, \emph{uniformly in $M$},
  \begin{equation}
    \label{eq:window-split}
    \Bigl\lvert \unifavg e(\beta k_n) \Bigr\rvert \leq \abs{P_i(\beta)} + \frac{2\delta^i}{N-M} .
  \end{equation}

  \emph{The limit of $P_\ell$.}  If $\norm{b^c \beta} \geq \eta$ for infinitely many $c$, then
  since $\rho$ is continuous, $1$-periodic and $\abs{\rho} < 1$ off $\ZZ$, there is $\kappa > 0$
  with $\abs{\rho(\theta)} \leq 1 - \kappa$ when $\norm{\theta} \geq \eta$; infinitely many
  factors in \eqref{eq:Pell} are then $\leq 1 - \kappa$ and $P_\ell(\beta) \to 0$.
  Otherwise $\norm{b^c\beta} \to 0$; choosing $c_0$ with $\norm{b^c\beta} < 1/(2b)$ for
  $c \geq c_0$, the relation $\norm{b^{c+1}\beta} = b\norm{b^c\beta}$ (valid when
  $0 < \norm{b^c \beta} < 1/(2b)$) forces $\norm{b^{c_0}\beta} = 0$, i.e.\ $\beta \in E_0$.

  \emph{Conclusion.}  For $\beta \notin E_0$ and $\epsilon > 0$ fix $i$ with
  $\abs{P_i(\beta)} < \epsilon$; by \eqref{eq:window-split},
  $\Phi_L(\beta) \leq \epsilon + 2\delta^i/L < 2\epsilon$ once $L > 2\delta^i/\epsilon$.
\end{proof}

\begin{lemma}[Uniform spectral lemma]
  \label{lem:uniform-spectral}
  \lean{CantorSzemeredi.CantorData.PhiCircle, CantorSzemeredi.CantorData.measurable_Phi,
    CantorSzemeredi.CantorData.excSetSmul_countable,
    CantorSzemeredi.CantorData.tendsto_integral_Phi_zero}
  \leanok
  \uses{lem:irrational, def:Phi}
  Let $\sigma$ be a finite positive Borel measure on $\TT$ without atoms and $s \geq 1$ an
  integer.  Then $\displaystyle \lim_{L \to \infty} \int_{\TT} \Phi_L(s\alpha) \, d\sigma(\alpha) = 0$.
\end{lemma}

\begin{remark}[Lean design note]
  $\TT$ is \texttt{AddCircle (1 : \ensuremath{\mathbb{R}})}, and ``without atoms'' is the
  class \texttt{NullSingletonClass}.  Since the $k_n$ are integers, $\Phi_L$ is $1$-periodic
  and descends to $\TT$ as \texttt{CantorData.PhiCircle}; the dilation $\alpha \mapsto s\alpha$
  is the $\NNz$-scalar multiple $s \cdot \alpha$ on $\TT$.
\end{remark}

\begin{proof}
  \leanok
  \uses{lem:irrational}
  Each $\Phi_L$ is a countable supremum of continuous functions, hence Borel, and
  $0 \leq \Phi_L \leq 1$.  By Lemma~\ref{lem:irrational}, $\Phi_L(s\alpha) \to 0$ unless
  $s\alpha \in E_0$; the set of such $\alpha$ is countable, hence $\sigma$-null.  Dominated
  convergence.
\end{proof}

\begin{corollary}[Uniform ergodic theorem along $\Cantor$]
  \label{cor:uniform-ergodic}
  \lean{CantorSzemeredi.tendsto_uniformAvg_koopman_zero}
  \uses{lem:uniform-spectral, def:spectral-measure}
  If $f \in L^2(X)$ has atomless spectral measure $\sigma_f$ and $s \geq 1$, then
  $\bigl\lVert \unifavg T^{s k_n} f \bigr\rVert_2 \to 0$ as $N - M \to \infty$.
\end{corollary}

\begin{proof}
  \uses{lem:uniform-spectral, def:spectral-measure}
  By the spectral theorem,
  $\bigl\lVert \unifavg T^{s k_n} f \bigr\rVert_2^2
   = \int_{\TT} \bigl\lvert \unifavg e(s\alpha k_n) \bigr\rvert^2 d\sigma_f(\alpha)
   \leq \int_{\TT} \Phi_{N-M}(s\alpha)\, d\sigma_f(\alpha)$,
  which tends to $0$ by Lemma~\ref{lem:uniform-spectral}.
\end{proof}
